\assignq{4}{1}{Hartree--Fock: direct and exchange Coulomb matrix elements}{1}
The task refers specifically to Johnson's closed-shell
results, Eqs.~(3.64) and (3.66). Here we derive both
from the multipole expansion, including the sum
over the \emph{occupied} magnetic and spin projections
of shell \(b=(n_bl_b)\).

\subsection*{1. Radial, orbital and spin separation}
In atomic units, let a spin-orbital be
\[
\psi_a(\mathbf r,\sigma)=\frac{P_a(r)}r
Y_{l_am_a}(\Omega)\chi_{\mu_a}(\sigma),
\qquad\int_0^\infty P_a(r)^2\diff r=1.
\]
For four spin orbitals define
\[
g_{abcd}=\int\!\!\diff x_1\int\!\!\diff x_2\
\frac{\psi_a^*(x_1)\psi_b^*(x_2)
\psi_c(x_1)\psi_d(x_2)}{r_{12}}.
\]
The Coulomb operator has the spherical expansion
\[
\frac1{r_{12}}
=\sum_{k=0}^\infty\sum_{q=-k}^{k}
(-1)^q\frac{r_<^k}{r_>^{k+1}}\,
C^k_{-q}(\Omega_1)C^k_q(\Omega_2),
\]
where \(C^k_q=\sqrt{4\pi/(2k+1)}Y_{kq}\).
Define the \emph{Slater radial integral}
\[
\boxed{\begin{aligned}
R^k(a,b,c,d)
=\int_0^\infty\!\diff r_1\,P_a(r_1)P_c(r_1)
\int_0^\infty\!\diff r_2\,P_b(r_2)P_d(r_2)
\frac{r_<^k}{r_>^{k+1}}.
\end{aligned}}
\]
The \(r^2\) Jacobian is already cancelled
by the \(1/r\) factors in each radial orbital.

\subsection*{2. General angular matrix element}
Introduce
\(A^k_q(a,c)=
\langle l_am_a|C^k_q|l_cm_c\rangle\).
Then the multipole expansion separates as
\[
g_{abcd}=\delta_{\mu_a\mu_c}
\delta_{\mu_b\mu_d}
\sum_{kq}(-1)^q R^k(a,b,c,d)
A^k_{-q}(a,c)A^k_q(b,d).
\]
The Wigner--Eckart theorem gives
\[
\langle l_bm_b|C^k_q|l_am_a\rangle
=(-1)^{l_b-m_b}
\begin{pmatrix}l_b&k&l_a\\-m_b&q&m_a\end{pmatrix}
\langle l_b\Vert C^k\Vert l_a\rangle,
\]
where
\[
\boxed{\langle l_b\Vert C^k\Vert l_a\rangle
=(-1)^{l_b}\sqrt{(2l_b+1)(2l_a+1)}
\begin{pmatrix}l_b&k&l_a\\0&0&0\end{pmatrix}.}
\]
The rank \(k\) must satisfy
\(|l_a-l_b|\le k\le l_a+l_b\) and
\(l_a+l_b+k\) even.

\subsection*{3. The direct integral, \(g_{abab}\)}
In the direct product, each electron keeps
its own orbital and spin:
\[
g_{abab}=
\langle a(1)b(2)|r_{12}^{-1}|a(1)b(2)\rangle.
\]
Sum over the two spin values and all
\(2l_b+1\) magnetic projections in the
\emph{closed} \(b\) subshell. The
spherical-harmonic addition theorem gives
\[
\sum_{m_b=-l_b}^{l_b}
|Y_{l_bm_b}(\Omega)|^2
=\frac{2l_b+1}{4\pi}.
\]
This summed density is spherically symmetric,
so all multipoles \(k>0\) vanish. Only
\(k=0\), for which \(C^0_0=1\), contributes:
\[
\boxed{\sum_{m_b,\mu_b}g_{abab}
=2(2l_b+1)R^0(a,b,a,b).}
\tag{Johnson 3.64}
\]
The factor \(2\) comes from summing the two
possible spin projections of occupied shell \(b\).

\subsection*{4. The exchange integral, \(g_{abba}\)}
The exchange matrix element swaps the orbital
labels on the ket:
\[
g_{abba}=
\langle a(1)b(2)|r_{12}^{-1}|b(1)a(2)\rangle.
\]
Its spin factor is
\(\delta_{\mu_a\mu_b}\delta_{\mu_b\mu_a}\),
so only \emph{same-spin} electrons contribute.
The sum over \(\mu_b\) gives \emph{one}, not two.
The angular structure contains the squared
transition matrix elements
\(|\langle l_bm_b|C^k_q|l_am_a\rangle|^2\).
For fixed \(m_a\), 3-j orthogonality states
\[
\sum_{m_b,q}
\begin{pmatrix}l_b&k&l_a\\-m_b&q&m_a\end{pmatrix}^{\!2}
=\frac1{2l_a+1}.
\]
Consequently,
\[
\boxed{\sum_{m_b,\mu_b}g_{abba}
=\sum_{k}\frac{|\langle l_b\Vert C^k\Vert l_a\rangle|^2}
{2l_a+1}\,R^k(a,b,b,a).}
\tag{Johnson 3.66}
\]
Only ranks allowed by the triangle and
parity selection rules contribute.

\subsection*{5. Why these yield the Hartree--Fock equations}
The Coulomb expectation of an antisymmetrised
Slater determinant contains the combination
\[
\boxed{E_{\rm ee}=\frac12\sum_{a,b\ {\rm occupied}}
\bigl(g_{abab}-g_{abba}\bigr).}
\]
The first term is the direct, local average
electron--electron repulsion. The second is
a nonlocal exchange contribution arising
from the \emph{antisymmetry} of electron
wavefunctions, not a separate classical force.
Variational minimisation with orthonormal
orbitals gives Hartree--Fock equations with
direct and exchange operators. For a single
doubly occupied \(1s\) subshell, the direct
contribution counts both spin occupancies
but exchange removes the self-interaction
of the occupied spin-orbital.

\methodnote{The two target expressions are
\emph{shell-summed} matrix elements.
Without this sum, neither the factor
\(2(2l_b+1)\) nor the reduction to the
single \(k=0\) direct term follows.
Source: Walter R. Johnson,
\emph{Atomic Structure Theory}, chapter 3,
Eqs.~(3.55)--(3.66).}
